% Clase 4 --- El producto en la recta: construcción por el teorema de Tales
% (§8.2). Recta principal R con O, I=1, Y=y, X=x, Z=xy (valores ilustrativos
% x=2, y=1.5, xy=3, elegidos solo para que la figura sea legible). Recta
% secante R' por O con la "unidad duplicada" I' a distancia arbitraria
% a=1.3 y G' a distancia a·y (Tales, paso 1). El paso 2 repite Tales entre
% I'-X y G'-Z para obtener Z=xy sobre R. Construcción estándar de Tales
% para el producto (la transcripción da los mismos puntos O,I',G',Z, pero
% sin el detalle numérico de la razón).
\begin{center}
\begin{tikzpicture}[scale=1.15]

  % ---- parámetros de la construcción (coordenadas calculadas, no a ojo) ----
  % ang=35°, a=OI'=1.3 (unidad duplicada arbitraria); x=OX=2, y=OY=1.5.
  % Por Tales: OG'=a*y=1.95, Z=OX*OG'/OI'=x*y=3. Coordenadas de R' obtenidas
  % de (r*cos(ang), r*sin(ang)) con ang=35°: cos35=0.81915, sin35=0.57358.

  % ---- recta principal R (el eje numérico) ----
  \draw[figaxis] (-0.5,0) -- (3.7,0) node[right, figlbl] {$R$};

  % ---- recta secante R' ----
  \draw[figline] (-0.2048,-0.1434) -- (1.8840,1.3192)
        node[above right, figlbl] {$R'$};

  % ---- puntos sobre R ----
  \node[figptocerrado] (O) at (0,0)   {};
  \node[figptocerrado] (I) at (1,0)   {};
  \node[figptocerrado] (Y) at (1.5,0) {};
  \node[figptocerrado] (X) at (2,0)   {};
  \node[circle, inner sep=0pt, minimum size=5pt, draw=figred, line width=0.9pt, fill=figred] (Z) at (3,0) {};

  \node[figlbl, below=3pt] at (O) {$O$};
  \node[figlbl, below=3pt] at (I) {$I$};
  \node[figlbl, below=3pt] at (Y) {$Y$};
  \node[figlbl, below=3pt] at (X) {$X$};
  \node[figlbl, figred, below=3pt] at (Z) {$X{\cdot}Y$};

  % ---- puntos sobre R': I' a distancia 1.3, G' a distancia 1.5*1.3=1.95 ----
  \node[figptocerrado] (Ip) at (1.0649,0.7456) {};
  \node[figptocerrado] (Gp) at (1.5973,1.1185) {};

  \node[figlbl, right=2pt] at (Ip) {$I'$};
  \node[figlbl, right=2pt] at (Gp) {$G'$};

  % ---- paso 1 de Tales: I-I' y su paralela por Y hasta G' ----
  \draw[figline] (I) -- (Ip);
  \draw[figline] (Y) -- (Gp);

  % ---- paso 2 de Tales: I'-X y su paralela por G' hasta Z ----
  \draw[figred, line width=0.9pt] (Ip) -- (X);
  \draw[figred, line width=0.9pt] (Gp) -- (Z);

\end{tikzpicture}\\[0.5em]
{\small\itshape Construcción estándar de $x\cdot y$ por el teorema de Tales
(triángulos semejantes). \textbf{Paso 1} (azul): se traza $I'$ arbitrario
sobre la secante $R'$ y la paralela a $II'$ por $Y$, que determina $G'$ con
$OG'/OI' = OY/OI$. \textbf{Paso 2} (rojo): se traza la paralela a $I'X$ por
$G'$, que determina $Z$ sobre $R$ con $OZ/OX = OG'/OI'$; combinando ambas
razones, $OZ = OX \cdot OY / OI = x\cdot y$ (con $OI=1$). A diferencia de la
suma, el producto \textbf{sí} necesita fijar la unidad $I$.}
\end{center}
